problem:  
Let \( (M^n,g) \) be a compact, simply connected Riemannian manifold with **quasipositive curvature**, that is, \(\sec \ge 0\) everywhere and \(\sec >0\) somewhere. Assume that \(M\) admits an **effective isometric torus action** \(T^r \curvearrowright M\) with rank  
\[
r = \left\lfloor \frac{n}{2} \right\rfloor - 1,
\]
so the symmetry rank is one less than the maximal possible for positive curvature.  

In the **strictly positive curvature** case, existing symmetry‑rank results (e.g. due to Grove–Searle, Wilking, Fang–Rong, and others) impose strong **topological and homotopical restrictions** on such manifolds—often forcing their **rational cohomology** or **tangential homotopy type** to resemble that of a compact rank‑one symmetric space (CROSS), though a complete classification is not known.  

**Question:**  
Does there exist a compact, simply connected, quasipositively curved manifold admitting an effective torus action of symmetry rank \( r = \lfloor \tfrac{n}{2} \rfloor - 1 \) whose **rational cohomology ring** is not isomorphic to that of any CROSS?  
Equivalently, can weakening curvature from positive to quasipositive enlarge the universe of rational cohomology types among near‑maximal symmetry manifolds?  

Why is it a "good" problem:  
- **Profound insight:**  
  This problem cuts precisely at the **rigidity boundary** between positive and quasipositive curvature. It asks whether the smallest possible relaxation of symmetry—from maximal to “one less than maximal”—and of curvature—from everywhere positive to merely quasipositive—still rigidly enforces the algebraic topology typical of positively curved spaces. This probes whether curvature positivity is essential for the known high‑symmetry rigidity, or whether quasipositivity suffices.  

- **Bridge between fields:**  
  An answer likely requires interfacing **Riemannian geometry** (sectional curvature and torus actions), **Alexandrov geometry** (orbit‑space structure with lower curvature bounds), **rational homotopy theory** (classification of Betti numbers and homogeneous models), and potentially **Ricci flow or collapse theory** (for degeneration analysis toward nonnegatively curved limits). It therefore forges a deep link between geometric analysis, topology, and transformation group theory.  

- **Extreme simplicity and depth:**  
  Despite its succinct statement—just one rank and curvature step away from well‑understood territory—the problem is technically and conceptually formidable. It encapsulates the question: *Does quasipositive curvature truly introduce new topology under almost maximal symmetry?* Resolving this would expose how fragile geometric rigidity is with respect to curvature assumptions.  

- **Potential impact:**  
  A construction of such a manifold would reveal brand‑new examples of quasipositively curved spaces beyond classical models, altering expectations in the field. A proof of nonexistence would unify rigidity phenomena across both curvature regimes. Either way, the result would significantly advance understanding of how curvature and symmetry interact to constrain manifold topology.  

Thus the problem’s **minimal hypotheses**, its position at a **critical transition point in the rigidity range**, and its potential to alter our view of the geometry–topology–symmetry relationship make it a quintessentially **good research problem**.